% !TeX program = lualatex
% Compile twice: lualatex 231.tex
% All references and prose are embedded; no BibTeX database or external inputs.
\documentclass[11pt,a4paper]{article}
\usepackage[margin=24mm,headheight=14pt]{geometry}
\usepackage{fontspec}
\setmainfont{Latin Modern Roman}
\setsansfont{Latin Modern Sans}
\setmonofont{Latin Modern Mono}
\usepackage{amsmath,amssymb,mathtools}
\usepackage{unicode-math}
\setmathfont{Latin Modern Math}
\usepackage{microtype}
\usepackage{enumitem,xurl,needspace}
\usepackage[dvipsnames]{xcolor}
\usepackage{hyperref}
\hypersetup{unicode=true,colorlinks=true,linkcolor=MidnightBlue,urlcolor=MidnightBlue,
 pdftitle={Research Notebook: Section 231},pdfauthor={Research notebook},bookmarksnumbered=true}
\usepackage{bookmark}
\usepackage{tocloft}
\setlength{\cftsecnumwidth}{3em}
\cftsetpnumwidth{2.5em}
\cftsetrmarg{4em}
\usepackage{titlesec}
\titleformat{\section}{\Large\bfseries\raggedright}{\thesection}{0.7em}{}
\usepackage{fancyhdr}
\pagestyle{fancy}\fancyhf{}
\fancyhead[L]{\small\sffamily Open Problems | Research notebook}
\fancyhead[R]{\small 231 | 27 September 2026}
\fancyfoot[C]{\thepage}
\setlength{\parindent}{0pt}
\setlength{\parskip}{5pt plus 1pt}
\setlength{\emergencystretch}{3em}
\setcounter{tocdepth}{1}
\setcounter{secnumdepth}{1}
\setlist[itemize]{leftmargin=3em,itemsep=5pt,topsep=4pt,parsep=0pt}
\allowdisplaybreaks
\newcommand{\N}{\mathbb{N}}
\newcommand{\Z}{\mathbb{Z}}
\newcommand{\Q}{\mathbb{Q}}
\newcommand{\R}{\mathbb{R}}
\newcommand{\C}{\mathbb{C}}
\newcommand{\F}{\mathbb{F}}
\newcommand{\A}{\mathbb{A}}
\newcommand{\PP}{\mathbb P}
\newcommand{\E}{\mathbb{E}}
\newcommand{\eps}{\varepsilon}
\newcommand{\dd}{\,\mathrm d}
\newcommand{\sref}[2]{\hyperlink{source:#1:#2}{[S#2]}}
\newcommand{\eref}[2]{\hyperlink{extra:#1:#2}{[E#2]}}
% Turn the next switch off for a shorter reading copy. The quotations preserve
% every exactTarget field, including qualifications omitted from a short summary.
\newif\ifshowcatalogue
\showcataloguetrue
\newcommand{\cataloguescope}[1]{\ifshowcatalogue
 \paragraph{Exact scope in the supplied catalogue.}
 \begin{quote}\small #1\end{quote}\fi}

% Independently compilable extraction; original section text retained.
\numberwithin{equation}{section}
\begin{document}
\setcounter{section}{230}
% ================================================================
% problemId: problem.zauners-conjecture-on-the-existence-of-sic-povms-in-all-dimensions
% source releaseRank: 231; source releaseStatus: open
% exactTarget SHA-256: 68bf69ecacd2c9a9f9bc585e4a748a5d3c2a8bee7e1c5fbe31ee77cd639bca4d
\clearpage
\section{\texorpdfstring{Zauner's Conjecture: Unrestricted Complex SIC Existence in Every Dimension}{Zauner's Conjecture: Unrestricted Complex SIC Existence in Every Dimension}}\label{231}
\subsection{Definitions and mathematical statement}
In $\C^d$ with its standard Hermitian inner product, seek $d^2$ unit vectors satisfying
\[
 |\langle v_i,v_j\rangle|^2=\frac1{d+1}\quad(i\ne j)
\]
for every integer $d\ge2$. Their lines form a symmetric informationally complete configuration. Writing $P_i=v_iv_i^*$, the same configuration has
$\sum_iP_i=dI_d$; hence the positive operators $M_i=P_i/d$ sum to $I_d$ and form a POVM, with
$\operatorname{Tr}(M_iM_j)=1/[d^2(d+1)]$ for $i\ne j$. The assertion asks existence only. Additional group covariance, an order-three symmetry, algebraic coordinates, or a uniform construction algorithm are not required.
\par\smallskip\textit{Formulation and concept references:} \sref{231}{1}.
\cataloguescope{Decide whether for EVERY integer d>=2 there exist d\textasciicircum{}2 unit vectors v\_1,...,v\_(d\textasciicircum{}2) in the standard complex Hilbert space C\textasciicircum{}d with |<v\_i,v\_j>|\textasciicircum{}2=1/(d+1) whenever i!=j. Their complex lines are distinct. Equivalently, their rank-one projectors P\_i=v\_i v\_i\textasciicircum{}* have sum\_i P\_i=d I\_d and define SIC-POVM elements M\_i=P\_i/d; for i!=j the normalized trace is Tr(M\_i M\_j)=1/[d\textasciicircum{}2(d+1)]. No Weyl-Heisenberg covariance, order-three symmetry, explicit coordinates or uniform construction algorithm is imposed.}
\subsection{Short English statement}
Does every complex dimension admit the maximal symmetric family of equiangular lines needed for a symmetric informationally complete quantum measurement?
% BEGIN RESEARCH ATTEMPT 231
\subsection{Research attempt}

\paragraph{Current frontier.}
The formulation source relates SICs to tight frames and complex
projective designs, and distinguishes numerical evidence from an
all-dimensional existence theorem. \sref{231}{1}.
A recent adjacent result of Cobucci and coauthors, revised in February
2026 and published in Physical Review Letters, shows that optimizing a
natural measure of non-projectivity does not in general select SICs
beyond qubits. That is a caution against changing the objective, not a
nonexistence theorem for SICs. \eref{231}{1}.
No all-dimensional existence or nonexistence result was verified in this
source check. The present target imposes no Weyl--Heisenberg covariance;
failure of a particular covariance ansatz would not disprove it.

\paragraph{Attack route.}
The first route minimizes a frame potential on a compact space. It needs
a proof that its lower bound is attained, not merely a proof of existence
of a minimizer. A second places a regular simplex in the real space of
traceless Hermitian matrices; it needs the vertices to lie on the
rank-one-projector submanifold. A third tensors known small SICs, which
preserves tightness but might not preserve equiangularity. We combine
these routes to derive an exact gap identity and construct positive-gap
stationary configurations. The aim is to test whether variational or
dimension-building arguments can close the existence gap without adding
a covariance restriction.

\paragraph{Attempt.}
\textbf{1. An exact nonnegative gap functional.}
Fix $d\geq2$, set $N=d^2$, and let $v_1,\ldots,v_N$ be arbitrary unit
vectors. Put $P_i=v_iv_i^*$, $s_{ij}=\operatorname{Tr}(P_iP_j)$ and
\[
 S=\sum_iP_i,\qquad
 \Phi_1=\sum_{i,j}s_{ij}=\operatorname{Tr}(S^2),\qquad
 \Phi_2=\sum_{i,j}s_{ij}^2.
\]
Since $\operatorname{Tr}S=N$, the sum-of-squares inequality on its
$d$ eigenvalues gives
\begin{equation}\label{231:tight}
 \Phi_1\geq N^2/d=d^3,
 \qquad \Phi_1=d^3\Longleftrightarrow S=dI_d.
\end{equation}
Write $c=(d+1)^{-1}$. Expanding the off-diagonal squares and using
$s_{ii}=1$ and $N=d^2$ gives the exact identity
\begin{equation}\label{231:gap}
 \Phi_2-\frac{2d^3}{d+1}
 =\sum_{i\ne j}(s_{ij}-c)^2
       +\frac2{d+1}\bigl(\Phi_1-d^3\bigr).
\end{equation}
For clarity, the constant in
$\sum_{i\ne j}(s_{ij}-c)^2=\Phi_2-2c\Phi_1$
vanishes because $-N+2cN+c^2N(N-1)=0$.
Both terms on the right of \eqref{231:gap} are nonnegative.
Equality holds exactly when all off-diagonal overlaps have the target
value and $S=dI_d$. Indeed, the overlaps alone also imply tightness:
then $\operatorname{Tr}(S^2)=N+N(N-1)c=d^3$.
Thus attaining this bound is exactly the unrestricted existence problem,
not a relaxation to a different quantum measurement.

The same lower bound follows by considering
$A=\sum_i|v_i\otimes v_i\rangle\langle v_i\otimes v_i|$ on the
symmetric square, whose dimension is $d(d+1)/2$:
$\operatorname{Tr}A=N$ and $\operatorname{Tr}A^2=\Phi_2$.
This is the projective-design interpretation used in \sref{231}{1}.
The domain $(S^{2d-1})^{d^2}$ is compact, so $\Phi_2$ has a global
minimum. Compactness does not prove that the minimum equals its
Cauchy--Schwarz lower bound. That unproved equality is exactly what a
variational solution still needs.

\textbf{2. The regular-simplex construction has an extra quadratic constraint.}
Pass to $Q_i=P_i-I_d/d$ in the $D=d^2-1$ dimensional real inner-product
space of traceless Hermitian matrices, with inner product
$\operatorname{Tr}(AB)$. A SIC would give
\begin{equation}\label{231:simplex}
 \operatorname{Tr}(Q_i^2)=\frac{d-1}{d},\qquad
 \operatorname{Tr}(Q_iQ_j)=-\frac1{d(d+1)}\quad(i\ne j),\qquad
 \sum_iQ_i=0.
\end{equation}
A regular simplex with those inner products exists abstractly in every
such real vector space. However, being the translate of a rank-one
projector also requires
\begin{equation}\label{231:rankone}
 Q_i^2=\left(1-\frac2d\right)Q_i+\frac{d-1}{d^2}I_d.
\end{equation}
Conversely, a traceless Hermitian $Q$ satisfying
\eqref{231:rankone} has eigenvalues $(d-1)/d$ and $-1/d$.
The trace forces multiplicity one for the former. Hence $Q+I_d/d$ is
indeed a positive rank-one projector. This proves that the omitted
constraint is both necessary and sufficient, not a cosmetic positivity
convention.

An orthogonal rotation of the real space preserves
\eqref{231:simplex}, but need not preserve \eqref{231:rankone}.
For example, apply the orthogonal map $Q\mapsto-Q$ to any actual SIC
in dimension $d>2$. The reconstructed matrix is
\[
 P'=\frac{2}{d}I_d-P,
\]
which has eigenvalue $2/d-1<0$ along the original vector. Thus the
reflected simplex satisfies all the proposed Gram conditions but is
not a family of positive projectors. The dimension-three example below
makes this an explicit test independent of any unproved large-dimensional
existence. The real-simplex construction cannot simply discard the
matrix quadratic equations.

\textbf{3. Exact seeds in dimensions two and three.}
Let $\omega=e^{2\pi i/3}$. In dimension two take
\[
 (1,0),\qquad
 \left(\frac1{\sqrt3},\sqrt{\frac23}\,\omega^k\right)
       \quad(k=0,1,2).
\]
For two different nonfirst vectors, the squared overlap is
$|1+2\omega^{k-\ell}|^2/9=1/3$; overlap with the first vector also
has square $1/3$. In dimension three take the nine vectors consisting
of the cyclic coordinate permutations of
\[
 \frac1{\sqrt2}(0,1,-\omega^k),\qquad k=0,1,2.
\]
Two vectors with the same zero coordinate have squared overlap
$|1+\omega^{k-\ell}|^2/4=1/4$ when distinct. With different zero
coordinates they have just one overlapping nonzero coordinate, so
again the square is $1/4$. These exact elementary seeds are well-known
small SICs, not claimed as new constructions. The companion calculation
checks their Gram matrices and tight-frame sums in exact algebraic
arithmetic.

\textbf{4. Tensoring gives a stationary obstruction, not induction.}
Suppose SICs are available in dimensions $a,b\geq2$. Their Cartesian
tensor products give $a^2b^2$ unit vectors in dimension $d=ab$ and a
tight frame, since
\[
 \sum_{i,j}P_i\otimes R_j=aI_a\otimes bI_b=dI_d.
\]
But overlaps between distinct vectors can be $1/(a+1)$, $1/(b+1)$,
or $1/[(a+1)(b+1)]$, depending on which factors agree. They are not all
$1/(ab+1)$. More quantitatively, factorization of the full ordered
frame-potential sum gives
\begin{align}
 \Phi_2^{\rm tensor}
    &=\frac{2a^3}{a+1}\frac{2b^3}{b+1}
      =\frac{4d^3}{(a+1)(b+1)},\notag\\
 \Phi_2^{\rm tensor}-\frac{2d^3}{d+1}
    &=\frac{2d^3(a-1)(b-1)}{(a+1)(b+1)(ab+1)}>0.
       \label{231:tensorgap}
\end{align}
Thus the most immediate dimension-building construction fails by a
specified positive amount.

It has an additional useful feature: it is stationary for $\Phi_2$ on
the full product of unit spheres, not just under tensor-product variations.
The first variation with respect to a vector $v_i$ is proportional to
\[
 \operatorname{Re}\left\langle\delta v_i,
             \Bigl(\sum_j s_{ij}P_j\Bigr)v_i\right\rangle.
\]
For a SIC in dimension $a$,
\[
 \sum_j s_{ij}P_j
 =P_i+\frac1{a+1}(aI_a-P_i)
 =\frac{a}{a+1}(I_a+P_i).
\]
It therefore sends $v_i$ to $2a/(a+1)$ times $v_i$. For the tensor
family the corresponding operator factorizes, sending $v_i\otimes w_k$
to
$4ab/[(a+1)(b+1)]$ times that same vector. Every real tangent variation
to the unit sphere is orthogonal in real inner product to its vector,
so the first variation vanishes, even for variations that are not
tensor-product variations.

Using the exact dimensions $2$ and $3$ gives an explicit stationary
configuration in dimension $6$ with
\begin{equation}\label{231:six}
 \Phi_2=72,\qquad \frac{2d^3}{d+1}=\frac{432}{7},\qquad
 \text{gap}=\frac{72}{7}.
\end{equation}
This refutes the proposed shortcut that all stationary configurations
must be SICs. No assertion that this configuration is a local minimum
is made: first-order stationarity does not determine its Hessian.
A descent proof would need a mechanism to escape such positive-gap
critical configurations and would still need to identify the global
minimum.

\paragraph{Stress test.}
The objective in \eqref{231:gap} is exactly tied to the prescribed
overlap condition. It is not the different non-projectivity objective
studied in \eref{231}{1}. The regular-simplex test retains all its
linear and Gram constraints; its failure occurs precisely at the
projector constraint. Tensor tightness and the count $d^2$ are checked,
but neither substitutes for equal pairwise overlaps.

All finite examples and the value \eqref{231:six} are exact algebraic
calculations, not small numerical residuals promoted to existence in
unbounded dimension. Conversely, a successful numerical search in every
dimension tested would not cover every integer $d$. No order-three
symmetry, group covariance, or algorithmic construction requirement has
been added to the conjecture. A failure confined to any such ansatz
would remain a failure of that method.

\paragraph{Outcome.}
\textbf{Outcome: Exploratory approach with unresolved gap.}

The attempt derives an exact sum-of-squares/tightness gap identity,
identifies the quadratic constraint omitted by the abstract-simplex
route, and constructs an explicit positive-gap stationary configuration
by tensoring exact small SICs. These are mathematical diagnostics of the
proposed existence methods, with no claim of novelty priority for their
elementary frame-theoretic ingredients.

What remains is a proof that the global minimum in
\eqref{231:gap} is zero in every dimension, or an obstruction showing a
positive minimum in at least one dimension. Neither conclusion is
obtained. The stationary tensor example is not a SIC nonexistence
example.
% END RESEARCH ATTEMPT 231
\subsection{Sources}
\begin{itemize}\raggedright
\item[\textnormal{[S1]}]\hypertarget{source:231:1}{} Joseph M. Renes, Robin Blume-Kohout, A. J. Scott, Carlton M. Caves, Symmetric Informationally Complete Quantum Measurements, arXivquant-ph/0310075v1,13October2003,8pages.
\url{https://arxiv.org/pdf/quant-ph/0310075v1}.
{\small\textit{Catalogue formulation source.}}
% BEGIN ADDED REFERENCES 231
\item[\textnormal{[E1]}]\hypertarget{extra:231:1}{}
G. Cobucci, R. Brinster, S. Khandelwal, H. Kampermann, D. Bru\ss,
N. Wyderka and A. Tavakoli, \emph{Maximally non-projective measurements
are not always symmetric informationally complete}, Physical Review
Letters \textbf{136} (2026), 060201; arXiv:2508.03652v2,
9 February 2026, Abstract and the non-projectivity optimization framework.
\url{https://arxiv.org/abs/2508.03652v2}.
Consulted 27 September 2026; used only to distinguish objectives, not as
an all-dimensional SIC existence or nonexistence theorem.
% END ADDED REFERENCES 231
\end{itemize}

\end{document}
